% \vspace{-10pt}
\section{Related Work}
% \vspace{-10pt}

\label{sec:related}
\paragraph{Discrete diffusion language models}
Discrete diffusion models adapt the iterative denoising paradigm to categorical spaces.
Early work on discrete diffusion studies how to design corruption kernels for categorical data, including uniform, structured, nearest-neighbor, and absorbing-state transitions~\citep{D3PM,multinomial_diffusion,ctdd}.
In language modeling, absorbing-state masking has become a dominant template, in which models learn to recover clean tokens from partially masked sequences~\citep{MDLM,MD4}.
Other formulations improve or reinterpret this paradigm through discrete score estimation, conditional-distribution modeling, or importance-sampling-based refinement of the reverse process~\citep{SEDD,MDM_secret,edlm, DFM}.
Recent scaling efforts further show that diffusion language models can be trained or adapted at scaled model sizes up to 100B, substantially improving generation quality and downstream performance~\citep{DiffuLLaMA,llada-moe,LLaDA,dream7b,sdar,dmax,llada20,llada21}.
Despite this progress, most methods still define learning through a prescribed forward corruption process.
\method departs from this paradigm by learning refinement directly from sampler-induced intermediate states.
\vspace{-4pt}

\paragraph{Looped transformers}
Looped transformers reuse shared parameters across depth to iteratively refine a recurrent state, beginning with Universal Transformers~\citep{dehghani2019universal} and extending to recent reasoning-oriented and recurrent-depth language models~\citep{saunshi2025latentthoughts,geiping2025recurrentdepth}.
\method belongs to this family but uses the model-generated token draft itself as the recurrent state.
This draft-space design makes each loop directly supervisable by the clean target, allowing preceding loops to be detached and avoiding BPTT across refinement iterations.
LoopUS similarly uses random deep supervision to train detached latent refinement trajectories~\citep{park2026loopus}; \method differs by coupling looped computation with sampler-induced diffusion refinement and block-parallel draft generation.
\vspace{-4pt}

\paragraph{Refinement-based language generation}
Drafting and refinement have been studied as alternatives to purely left-to-right AR generation.
Edit-based generation methods revise sequences through insertion, deletion, or iterative rewriting, providing early evidence that language generation can be organized as refinement rather than strict autoregression~\citep{levenshtein,diffuser}.
Corrector-style samplers further model refinement as a reverse causal process, in which tokens are corrected while conditioning on already generated future context~\citep{correctorar}.
In diffusion language models, several methods augment masked-diffusion backbones with token-level correction or remasking heads that decide which positions should be revisited during generation~\citep{informedcorrector,remedi,prism, GIDD}.
Other approaches model refinement at the sequence level, combine diffusion proposals with AR verification and rejection sampling~\citep{selfspec,tidar}, or employ self-correction~\citep{rediff,corrective-dlm,llada21}.
These methods show that refinement is a useful mechanism for improving draft quality and correcting local errors.
Nevertheless, most of them still operate within a predefined masking or corruption framework, so the refinement dynamics remain tied to hand-designed intermediate states.
By contrast, \method treats model-generated drafts themselves as the intermediate states and learns a forward-free refinement operator.
\vspace{-4pt}

\paragraph{Efficient decoding for language models}
Improving inference efficiency is a central challenge for language models, especially when generation requires many sequential forward passes.
For AR models, a large body of work accelerates decoding through multi-token prediction, speculative drafting, verification, and lightweight draft models~\citep{mtp-parallel,deepseek-v3,eagle1,eagle2,eagle3}.
Diffusion language models offer a different opportunity: they can update multiple tokens per forward pass, but this parallelism often introduces a quality--speed trade-off.
Prior work improves diffusion decoding by controlling confidence thresholds~\citep{diff-speed-conf, fast-dllm}, prediction entropy~\citep{eb-sampler}, or sampling policies~\citep{seed-diffusion, trace, credit-decoding}.
Other methods transfer diffusion-forcing ideas to discrete generation~\citep{diffusion-forcing,D2F}, approximate KV-cache reuse~\citep{fast-dllm,fast-dllm2, dKV_Cache}, or combine diffusion proposals with autoregressive likelihood correction~\citep{selfspec,tidar, dflash}.
These approaches primarily accelerate or stabilize inference under an existing diffusion or hybrid decoding process.
\method further improves the underlying refinement formulation itself, aligning training with the drafts used at inference time.
